%%%%%%%%%%%%%%%%%%%%%%%%%%%%%%%%%%%%%%%%%%%%%%%%%%%%%%%%%%%%%%%%%%%%%%%%
\section{What would improve the bounds}\label{sec:improve}
%%%%%%%%%%%%%%%%%%%%%%%%%%%%%%%%%%%%%%%%%%%%%%%%%%%%%%%%%%%%%%%%%%%%%%%%

Theorem~\ref{thm:main} leaves the loss $L(n)$ between about $\log n$ and $C\sqrt n$. This section
states what would lower the upper estimate $C\sqrt n$, why some natural routes cannot, and what
weaker targets remain. The propositions and the argument that gluing cannot help have complete
written proofs, which have not been checked independently; the expected value of $\beta$ is a
heuristic.

\subsection{Short separators for a fraction of meanders}

The proof of the lower bound applies one separator with about $\sqrt N$ ports to every meander. A
separator with fewer ports may exist only for some meanders; then the argument must also pay for how
rare these meanders are.

\begin{proposition}[fraction criterion]\label{prop:fraction}
Let $0<\tau<1$ and $c_{\mathrm f}>0$. Suppose that for all large $N$ a fraction at least
$e^{-c_{\mathrm f}N^\tau}$ of the meanders of order $N$ have a separator with at most
$c_{\mathrm f}N^\tau$ ports whose two disks each contain between $N/2$ and $3N/2$ of the $2N$
crossings. Then $L(n)\le K_{\mathrm f}\,n^\tau$ for all $n$, for some constant $K_{\mathrm f}$.
\end{proposition}

\begin{proof}[Key idea]
Run the proof of Proposition~\ref{prop:reverse} on these meanders only. Their number is at least
$e^{-c_{\mathrm f}N^\tau}M_N$ and at most a polynomial times $e^{2c_\Omega c_{\mathrm f}N^\tau}\max M_aM_b$,
so $L(a)+L(b)\le L(N)+B'N^\tau$. The descent of Section~\ref{sec:descent} goes through with
$N^{\tau-1}$ in place of $N^{-1/2}$, because the geometric sums still converge.
\end{proof}

With $\tau=1/2$ and fraction one this is the lower bound of Theorem~\ref{thm:main}. The criterion
allows separators of any shape. For example, if every face of the map can be reached from one fixed
face by crossing at most $\omega$ edges, the map has a balanced separator with at most $2\omega+1$ ports.

\subsection{Shallow meanders}

The simplest balanced cut is a vertical line through the middle gap of the river. It always exists,
it splits the crossings evenly, and its ports are the arches that pass over the gap, so its cost
depends on how many arches do. Say that an arch system has \emph{depth} at most $d$ if no point lies
under more than $d$ of its arches, and let $M^{(d)}_n$ count meanders of order $n$ whose upper and
lower arch systems both have depth at most $d$. We call these \emph{shallow} meanders.

\begin{proposition}[depth criterion]\label{prop:depth}
Suppose that for some $K_{\mathrm d}$ and $\beta>0$
\[ M^{(d)}_n\;\ge\;\mu^n\exp\!\big(-K_{\mathrm d}\,n\,d^{-\beta}-K_{\mathrm d}\,d\big)\qquad\text{for all }n,d\ge1 .\]
Then $L(m)\le K_{\mathrm d}'\,m^{1/(1+\beta)}$ for all $m$, for some constant $K_{\mathrm d}'$.
\end{proposition}

\begin{proof}
In a meander of depth at most $d$, at most $2d$ arches pass over the middle gap. By
Proposition~\ref{prop:linrec}, $M^{(d)}_{2m}\le e^{2B_{\mathrm{lin}}d}M_m^2$ with $B_{\mathrm{lin}}=\log(64R_\star^3)$.
With the hypothesis at $n=2m$,
\[ 2m\log\mu-2K_{\mathrm d}md^{-\beta}-K_{\mathrm d}d\;\le\;2B_{\mathrm{lin}}d+2\log M_m , \]
so $L(m)\le K_{\mathrm d}\,m\,d^{-\beta}+(B_{\mathrm{lin}}+K_{\mathrm d}/2)\,d$.
Take $d=\lceil m^{1/(1+\beta)}\rceil$.
\end{proof}

The exponent $1/(1+\beta)$ is below $1/2$ exactly when $\beta>1$. So shallow meanders improve on the
square root if and only if confining a meander to depth $d$ costs less than $d^{-1}$ per crossing in
the exponent.

\emph{Gluing short meanders never helps.} Let $\mu_d=\lim_n(M^{(d)}_n)^{1/n}$ be the growth rate of
shallow meanders. The only unconditional lower bounds for $M^{(d)}_n$ come from building shallow
meanders out of short ones, by products. A meander of order $n_1\le d$ has depth at most $n_1$, and a
product of many of them, combined in a balanced way, adds only a few levels of nesting. This gives
$\log(\mu/\mu_d)\lesssim L(n_1)/n_1$ with $n_1\approx d$. If $L(n_1)=O(n_1^{\tau})$, that is $\beta=1-\tau$, and Proposition~\ref{prop:depth} returns the exponent
$1/(2-\tau)$. Since
\[ \frac1{2-\tau}-\tau\;=\;\frac{(1-\tau)^2}{2-\tau}\;\ge\;0, \]
the exponent never improves; from $\tau=1/2$ it gives $2/3$, worse than the square root. A proof along
these lines must produce long, thin meanders that are not concatenations of short ones.

\emph{What is expected.} Heuristically, $\beta\approx1.92$. Typical arch heights in large meanders, the
numbers of arches nested above a point, are predicted by a KPZ-type argument to grow like $n^\nu$ with $\nu\approx0.52$, and Monte Carlo measurements on the
closely related semi-meanders, where the river is a half-line, give $0.518$~\cite{Golinelli}. For a process whose height after $n$ steps is about
$n^\nu$, staying below height $d$ for a long time costs a factor of about
$\exp(-\mathrm{const}\cdot d^{-1/\nu})$ per step, which suggests
$\log(\mu/\mu_d)\approx\mathrm{const}\cdot d^{-1/\nu}$ and $\beta=1/\nu\approx1.92$. If so, and if the
hypothesis of Proposition~\ref{prop:depth} holds with this $\beta$, it gives $L(n)=O(n^{0.34})$. For comparison, pairs of arch systems without the
single-loop condition have $\beta=2$ exactly: those of depth at most $d$ grow like
$\sigma_d=(2\cos(\pi/(d+2)))^4$ per order, and $16-\sigma_d\approx32\pi^2/(d+2)^2$.

\emph{A cleaner sufficient condition.} The hypothesis of Proposition~\ref{prop:depth} would follow, with
$\beta$ up to $2$, from a correlation statement: that
$\mu_d/\sigma_d\ge(\mu/16)\exp(-O(d^{-\beta'}))$ for some $\beta'>1$, with uniform prefactors. In words, among pairs of arch systems of depth at most
$d$, the probability of forming a single loop decays at nearly the same exponential rate as for all
pairs. The growth rates $\mu_d$ can be computed exactly for small $d$ by transfer matrices, finite
matrices that build a shallow meander one crossing at a time, so this can be tested before a proof is attempted.

\subsection{Existence of an exponent is not easier}

One might hope to prove that some exponent exists, $u_n=n^{-\alpha+o(1)}$ for some finite $\alpha$,
without finding it. This is at least as hard as a lower bound on $M_n$ of the form $\mu^n$ times a
power of $n$: if such an $\alpha$ exists,
then $L(n)\le(\alpha+1)\log n$ for large $n$, that is $M_n\ge\mu^nn^{-\alpha-1}$. By
Corollary~\ref{cor:halving}, such a polynomial lower bound is the same as $D(m)\le\chi\log m+O(1)$ for
some $\chi$, that is, $M_{2m}=O(m^{\chi}M_m^2)$. No method in this paper gives that.

\subsection{A weaker target: subexponentiality}\label{sec:subexp}

A summable positive sequence $(a_n)$ is called \emph{subexponential} if
\[ \frac1{a_N}\sum_{m=1}^{N-1}a_ma_{N-m}\;\longrightarrow\;2\sum_na_n \qquad(N\to\infty). \]
In words, a large total is almost always made of one large part and one small part. Take $a_n=u_n$. By
\eqref{eq:cutprob} the left side is exactly $\mathbb E_N[\#\text{cuts}]$. So
\[ (u_n)\text{ is subexponential}\quad\Longleftrightarrow\quad\mathbb E_N[\#\text{cuts}]\to2U . \]
Since $\Pr_N(\text{cut at }m)\to u_m$ for each fixed $m$, by Theorem~\ref{thm:ratio}, the cuts near
the two ends already contribute $2U$ in the limit. So subexponentiality says that cuts with both sides
large carry vanishing total probability: for every $\epsilon>0$ there is $k_0$ with
$\sum_{k_0\le k\le N-k_0}\Pr_N(\text{cut at }k)\le\epsilon$ for all large $N$. Intuitively, when a large
meander has a cut, one side is almost always small.

This is a different target from a polynomial lower bound, and neither is known to imply the other. It
allows fast decay: sequences like $e^{-\sqrt n}$ are subexponential. It would have a concrete consequence. Since every meander is a sequence of primes and
$\sum u_n$ is finite, the standard theory of subexponential sequences~\cite{CNW} would give a positive limiting
fraction of prime meanders, $P_n/M_n\to(1-\mathcal P(1/\mu))^2$; the exact hypotheses of that theory
remain to be checked here. The exact data are consistent with it: $\mathbb E_N[\#\text{cuts}]$
decreases, to $0.245$ at $N=32$, while $2U\ge0.2064$ and the estimate of $U$ gives $2U\approx0.207$;
the middle part with $k_0=5$ is $0.0127$ at $N=32$ (Table~\ref{tab:cuts}). Proposition~\ref{prop:cutbound}
bounds each middle term, but its bound is not summable. What is missing is an argument that sees many
cut positions at once; local smoothness of $u_n$ alone cannot give it.
